\documentclass{article}

\usepackage[utf8]{inputenc}
\usepackage[german]{babel}
\usepackage{listings}
\usepackage{xcolor}

\lstset{
    basicstyle=\ttfamily\small,
    breaklines=true,
    backgroundcolor=\color{gray!10}
}

\begin{document}
    \begin{center}
        \Large \textbf{Huntsman Challenge 4: Sicherheit und Compliance Management} \\ [6pt]
        \normalsize
        Autor: Maximilian Jakob Maag B.Sc. \\
        Version: 1.0
    \end{center}

    \section{Intro}
    Diese Challenge befasst sich mit der Sicherung und Compliance der verwalteten Infrastruktur.
    Sie werden Sicherheitsrichtlinien durchsetzen, die Infrastruktur überwachen und sicherstellen, dass alle Systeme den Compliance-Anforderungen entsprechen.

    \section{Aufgaben}

    \subsection{Aufgabe 1}
    Erstellen Sie ein Sicherheits-Hardening Playbook, das auf alle verwalteten Server angewandt wird.
    \\ \\
    Das Playbook sollte:
    \begin{itemize}
        \item SSH Key Based Authentication erzwingen (Passwort Authentifizierung deaktivieren)
        \item Firewall Regeln konfigurieren (nur notwendige Ports öffnen)
        \item fail2ban installieren und konfigurieren
        \item sudo Zugriff auf Basis von Gruppen regeln
        \item Unwichtige Services deaktivieren
        \item SELinux oder AppArmor aktivieren
    \end{itemize}

    \subsection{Aufgabe 2}
    Implementieren Sie ein User Access Management Playbook.
    \\ \\
    Das Playbook sollte:
    \begin{itemize}
        \item Autorisierte Benutzer auf den Systemen verwalten
        \item SSH Public Keys zentral verwalten
        \item sudo Konfigurationen durchsetzen
        \item Inaktive Accounts deaktivieren
        \item Audit Logs für alle Zugriffe generieren
    \end{itemize}

    \subsection{Aufgabe 3}
    Erstellen Sie ein Security Scanning und Vulnerability Assessment Playbook.
    \\ \\
    Das Playbook sollte:
    \begin{itemize}
        \item Regelmäßig Sicherheits-Scans durchführen (z.B. mit Nessus, OpenVAS)
        \item Package Vulnerabilities überprüfen
        \item Firewall Regeln auf Offenheiten prüfen
        \item SSL/TLS Konfigurationen validieren
        \item Befunde in einem zentralen System sammeln und melden
    \end{itemize}

    \subsection{Aufgabe 4}
    Implementieren Sie Compliance Checking Playbooks für verschiedene Standards:
    \begin{itemize}
        \item CIS Benchmarks
        \item GDPR Anforderungen (Datenschutz)
        \item PCI DSS (für Payment Systems)
        \item Custom Company Policies
    \end{itemize}

    \subsection{Aufgabe 5}
    Erstellen Sie ein Playbook zur Patch Management und Updates.
    \\ \\
    Das Playbook sollte:
    \begin{itemize}
        \item Sicherheits-Updates automatisch installieren
        \item Critical Patches sofort einspielen
        \item Regelmäßige OS Updates nach Testlauf durchführen
        \item System nach Updates neu starten (mit Wartungsfenster)
        \item Compatibility mit laufenden Services prüfen
    \end{itemize}

    \subsection{Aufgabe 6}
    Implementieren Sie ein Logging und Monitoring Playbook für Security Events.
    \\ \\
    Das Playbook sollte:
    \begin{itemize}
        \item Centralized Log Collection einrichten (z.B. ELK Stack, Splunk)
        \item Failed Login Attempts loggen und monitoren
        \item Privilege Escalation Events tracken
        \item Datei-Änderungen auf kritischen Systemen monitoren
        \item Real-time Alerts für verdächtige Aktivitäten generieren
    \end{itemize}

    \subsection{Aufgabe 7}
    Erstellen Sie ein Incident Response Playbook.
    \\ \\
    Das Playbook sollte:
    \begin{itemize}
        \item Verdächtige Server isolieren können
        \item Forensische Daten sammeln
        \item Logs für Incident Investigation aufbewahren
        \item Automatische Notification an Security Team
        \item Rollback Prozeduren für kompromittierte Systeme
    \end{itemize}

    \subsection{Aufgabe 8}
    Dokumentieren Sie Ihre Sicherheitspolicies und Compliance Anforderungen.
    \\ \\
    Die Dokumentation sollte enthalten:
    \begin{itemize}
        \item Sicherheitsrichtlinien für alle Server
        \item Compliance Anforderungen und deren Umsetzung
        \item Security Incident Response Prozess
        \item Regelmäßige Security Audit Schedule
    \end{itemize}

\end{document}
